\subsection{LATTICES AND ORDERS}

Let V be a Q-vector space. A lattice in V is a free Z-submodule V of V such that the Q-linear map \(\alpha_{\mathcal{V},V} : \mathcal{V} \otimes_{Z} Q \to V\) induced by the injection of V into V is bijective. When V is finite dimensional, this is the same as saying that V is a Z-submodule of finite type which generates the Q-vector space V (recall that a torsion-free Z-module of finite type is free by Algebra, Chap. VII, §4, no. 4, Cor. 2); moreover, in this case our definition is a special case of that of Commutative Algebra, Chap. VII, §4, no. 1, Def. 1 (loc. cit., Example 3). If W is a vector subspace of V, and V is a lattice in V, then \(V \cap W\) is a lattice in W.

If V is a Q-algebra, an order in V is a lattice V in the underlying vector space that is a Z-subalgebra of V; the map \(\alpha_{V,V}\) is then an isomorphism of Q-algebras. If V is a unital Q-algebra, a unital order in V is an order in V containing the unit element.

Assume that V is a Q-bigebra, with coproduct c and counit \(\gamma\) . If V is a lattice in the vector space V, the canonical map \(i: V \otimes_{Z} V \to V \otimes_{Q} V\) is injective; a biorder in V is a unital order V in the unital algebra V such that \(\gamma(\mathcal{V}) \subset \mathbf{Z}\) and \(c(\mathcal{V}) \subset i(\mathcal{V} \otimes_{\mathbf{Z}} \mathcal{V})\) ; the maps

\[
\gamma_ {\mathcal {V}} \colon \mathcal {V} \to \mathbf {Z} \text { and } c _ {\mathcal {V}} \colon \mathcal {V} \to \mathcal {V} \otimes_ {\mathbf {Z}} \mathcal {V}
\]

induced by \(\gamma\) and c give V the structure of a Z-bigebra, and the map \(\alpha_{\mathcal{V},V}\) is then an isomorphism of Q-bigebras.

